% -----------------------------------------------------------------------
\section{Conclusiones y Trabajo Futuro}
\label{sec:conclusiones}

Se presentó un repositorio reproducible con 15 implementaciones de
Búsqueda Tabú (tres niveles) sobre cinco problemas clásicos de optimización
combinatoria, con experimentación automatizada.

\subsection*{Puntos de interés a nivel de implementación}

\begin{enumerate}
  \item \textbf{La progresión de niveles reduce el GAP global cerca de la
    mitad.}
    La transición de TS-Básico a TS-Avanzado reduce el GAP global de
    \SI{6.50}{\percent} a \SI{3.40}{\percent}, siendo CVRP el problema
    más beneficiado (reducción del \SI{67}{\percent} relativo).

  \item \textbf{El primer salto de complejidad concentra la mayor parte de
    la mejora.}
    TS-Intermedio reduce el GAP en 2.50\,pp respecto a
    TS-Básico, y TS-Avanzado reduce adicionalmente 0.60\,pp respecto a TS-Intermedio gracias a los mecanismos de largo plazo.

  \item \textbf{La perturbación informada es crítica en TSP.}
    El reinicio \emph{double-bridge} (4-opt) en TS-Avanzado es responsable
    de la mayor mejora individual observada: \SI{2.21}{\percent} de
    reducción en \texttt{kroA100}.

  \item \textbf{En Flowshop los tres niveles rinden de forma similar.}
    La heurística constructiva NEH entrega una solución inicial de buena
    calidad y los tres niveles de TS la mejoran poco (GAP medio entre
    \SI{1.74}{\percent} y \SI{1.75}{\percent}); el rendimiento en este
    problema está limitado por la vecindad de swap adyacente.

  \item \textbf{Job Shop es el problema más difícil del conjunto.}
    Los GAPs elevados en \texttt{la21} y \texttt{ft10} confirman la
    dificultad inherente del problema y la necesidad de vecindades más
    ricas (e.g., movimientos en el grafo de disyunción~\cite{balas1998one}).
\end{enumerate}

\subsection*{Trabajo futuro}

\begin{itemize}
  \item Incorporar modelos exactos (MILP) para las instancias de menor
    tamaño y comparar con los cotas inferiores.
  \item Implementar una cuarta configuración con TS poblacional
    (búsqueda tabú de trayectorias paralelas).
  \item Añadir instancias de mayor escala (e.g., TSP con 500+ ciudades,
    CVRP con 500+ clientes) y medir el escalado del GAP.
  \item Extender el repositorio a otros problemas como el de coloración
    de grafos o la mochila multidimensional.
  \item Automatizar la generación de figuras desde los cuadernos Jupyter
    e integrarlas directamente al PDF del artículo mediante \texttt{latexmk}.
\end{itemize}

En conclusión, la arquitectura progresiva de TS aquí validada permite
comprender el impacto incremental de cada mecanismo metaheurístico, y el
código publicado facilita su reproducción y extensión.
